\uuid{yjK7}
\exo7id{6397}
\titre{exo7 6397}
\auteur{potyag}
\organisation{exo7}
\datecreate{2011-10-16}
\isIndication{false}
\isCorrection{false}
\chapitre{Isométrie euclidienne}
\sousChapitre{Isométrie euclidienne}
\cpge{mpsi-mp2i.structures-algebriques.b, mpsi-mp2i.sous-espaces-affines, ptsi.espaces-vectoriels.b, mp-mpi.endomorphismes-euclidien.c, pc.endomorphismes-euclidiens.a, psi.endomorphismes-euclidien.a, pt.isometries-euclidien.a, mp-mpi.endomorphismes-euclidien.d, pc.endomorphismes-euclidiens.c, psi.endomorphismes-euclidien.d, pt.isometries-euclidien.c}
\module{Algèbre et géométrie}
\niveau{L3}
\difficulte{}

\contenu{
\texte{
Soient $f$, $g$ deux isométries du plan affine euclidien $\Rr^2$.
}
\begin{enumerate}
    \item \question{Montrer que si $g$ et $f$ commutent alors
$g=f\circ g\circ f^{-1}$  et $f$ conserve (globalement) l'ensemble
des points fixes de $g.$}
    \item \question{Décrire les cas dans lesquels $f$ et $g$
commutent ($f\circ g =g\circ f$).}
    \item \question{En déduire une description des sous-groupes
commutatifs d'isométries.}
\end{enumerate}
}
